\newcommand{\casefigure}[4]{%
\begin{figure*}[!t]
\centering
\includegraphics[width=#1]{#2}
\par\smallskip
\textbf{#3} #4
\end{figure*}
}

\newcommand{\casetable}[3]{%
\begin{table*}[!t]
\centering
\textbf{#1} #2
\par\smallskip
#3
\end{table*}
}

\section{Case Study}

\subsection{原始数据}

本文采用吉林省区域风电场群数据构建案例研究数据集。该数据集由两部分组成：一部分为风电场站的实测功率数据，另一部分为覆盖研究区域的网格化气象预报数据。实测功率数据用于描述风电场群的历史运行状态，气象预报数据用于提供预测窗口内的外部气象驱动。与单场站预测不同，本文关注区域层面的集群总功率预测，预测目标为研究区域风电场群的区域集群总功率。

功率数据覆盖 2022--2023 年吉林省多个风电场站的运行记录。研究区域风电场群包含 58 个风电场站。所有功率数据均以北京时间为时间基准，并统一整理为整点时间序列。研究区域内的场站分布如图 5 所示。

\casefigure{0.54\textwidth}{../../paper_image_use/data/fig05_station_distribution.png}{图 5.}{吉林区域风电场站空间分布。}

气象数据采用 Pangu-Weather 生成的逐小时预报产品。本文截取吉林省及其周边区域的网格化气象场作为模型输入，使气象信息能够覆盖风电场群所在的主要空间范围。气象变量包括海平面气压、2\,m 气温、不同气压层上的位势高度、比湿和温度，以及 10\,m、1000\,hPa 和 925\,hPa 风场。为更直接表示风电相关风况，风矢量变量被转换为风速及方向分量，最终形成 17 个气象输入通道。本文后续将这些变量按风场相关变量和背景气象变量组织，其中风场信息是影响风电出力变化的关键外部因素。每个样本中的气象预报按照第 3.1 节定义的业务时间协议与功率观测和预测目标进行对齐。

训练集与评估集按年份隔离划分。训练集使用 2022 年数据，并采用滑动窗口方式构造训练样本，以提高训练样本规模并覆盖不同运行状态。验证集和测试集使用 2023 年数据，均采用严格日样本设置，即每天 00:00 BJT 构造一个 24 小时预测样本。具体划分时，每月前 10 天作为验证集，其余日期作为测试集，验证样本和测试样本在预测日上互不重叠。该划分方式避免了随机划分可能造成的时间信息泄漏，也使模型选择和最终测试保持独立。

\subsection{实验设置}

为验证所提出模型在吉林区域风电场群 24 步集群总功率预测任务中的有效性，本文在固定的数据划分上进行实验。训练集、验证集和测试集分别对应第 4.1 节中给出的时间划分。所有模型均在相同的测试集上进行评价，并重复运行 12 次，报告均值和标准差。

与风电功率预测研究中常用的评价方式一致，本文采用 MAE、RMSE 和 \(R^2\) 评价确定性预测误差。设 \(P_i\) 和 \(\hat{P}_i\) 分别表示第 \(i\) 个测试点的真实集群总功率和预测集群总功率，\(\bar{P}\) 表示测试集中真实功率的均值，\(n\) 表示参与评价的测试点总数。在整体 24 步评价中，\(n\) 为测试样本数与预测步长的乘积；在逐步长评价中，\(n\) 为该预测步长下的测试样本数。由于模型输出为归一化功率，令 \(S_y\) 为训练集目标功率归一化尺度，\(y_i=P_i/S_y\)，\(\hat{y}_i=\hat{P}_i/S_y\)。归一化误差指标定义为：

\begin{equation}
\mathrm{MAE}=\frac{1}{n}\sum_{i=1}^{n}|y_i-\hat{y}_i|,
\end{equation}

\begin{equation}
\mathrm{MSE}=\frac{1}{n}\sum_{i=1}^{n}(y_i-\hat{y}_i)^2,
\end{equation}

\begin{equation}
\mathrm{RMSE}=\sqrt{\frac{1}{n}\sum_{i=1}^{n}(y_i-\hat{y}_i)^2}.
\end{equation}

实际功率尺度下的误差定义为：

\begin{equation}
\mathrm{MAE(MW)}=\frac{1}{n}\sum_{i=1}^{n}|P_i-\hat{P}_i|,
\end{equation}

\begin{equation}
\mathrm{RMSE(MW)}=\sqrt{\frac{1}{n}\sum_{i=1}^{n}(P_i-\hat{P}_i)^2}.
\end{equation}

\(R^2\) 用于衡量预测结果对测试集功率变化的解释能力，定义为：

\begin{equation}
R^2=1-\frac{\sum_{i=1}^{n}(P_i-\hat{P}_i)^2}{\sum_{i=1}^{n}(P_i-\bar{P})^2}.
\end{equation}

除整体 24 步平均结果外，本文还进一步分析不同预测步长上的误差变化，以检验模型在预测窗口内部的稳定性。

\casetable{表 1.}{实验设置。}{\input{../../paper_table_use/result/table01_experimental_setup.tex}}

对比模型包括 Weather-only 结构以及多种网格气象编码基线，包括 CNN、3D CNN、ViT、Swin 和 Video Transformer。所有基线均使用相同的数据划分和评价指标。本文模型与 Weather-only 的差异在于，本文模型进一步引入近期功率上下文，并将其用于条件化气象记忆；因此二者的对比用于检验功率状态信息是否能为未来气象驱动预测提供补充。

\subsection{基准模型对比}

表 2 给出了不同模型在测试集上的整体预测结果。本文模型在所有主要指标上均取得最优结果，MAE 为 \(0.0767\pm0.0006\)，RMSE 为 \(0.1041\pm0.0009\)，\(R^2\) 为 \(0.8330\pm0.0029\)。在实际功率尺度下，本文模型的 MAE(MW) 和 RMSE(MW) 分别为 \(377.9\pm3.1\) MW 和 \(512.7\pm4.4\) MW。

\casetable{表 2.}{基准模型整体预测结果。}{\input{../../paper_table_use/result/table02_benchmark_results.tex}}

与 Weather-only 相比，本文模型的 MAE 从 0.0800 降至 0.0767，RMSE 从 0.1072 降至 0.1041，\(R^2\) 从 0.8230 提升至 0.8330。这说明在相同气象输入条件下，引入近期功率上下文能够降低整体预测误差。除 Weather-only 外，3D CNN 是表现最好的气象编码基线，其 MAE 和 RMSE 分别为 0.0845 和 0.1138，仍高于本文模型。CNN、ViT、Swin 和 Video Transformer 的误差进一步增大。

\casefigure{\textwidth}{../../paper_image_use/result/fig06_benchmark_comparison.png}{图 6.}{基准模型误差对比。}

图 6 以误差条形式展示了各模型的 MAE 和 RMSE。图中排序与表 2 的整体结果一致：本文模型误差最低，Weather-only 位列第二，其余气象编码基线误差较高。Weather-only 与本文模型的差距小于其他气象编码基线，表明气象预报本身已经提供了强预测信号；但在此基础上加入功率上下文后，模型仍能进一步改善测试集表现。

\subsection{消融实验}

表 3 给出了功率上下文、功率编码器和气象模块相关消融结果。所有消融模型均与本文完整模型在相同测试集上比较，表中 \(\Delta\)MAE 和 \(\Delta\)RMSE 表示相对于完整模型的误差增幅。

为便于解释表中 Variant 的含义，本文将消融设置说明如下。Weather-only 移除所有功率侧条件化，仅保留气象编码器和 Query Decoder；Random power 和 Zero power 保留功率条件通路，但在功率编码器接收最近 \(L_p\) 步功率后，分别将其替换为同形状随机值和零值；Basic GCN 将完整模型中的点级、片段级和全局多尺度图编码以及跨尺度注意力融合，替换为单个 GCN 对最近 \(L_p\) 步功率进行编码。

气象侧消融用于检验 Weather Encoder 的结构设计。No memory 关闭气象记忆的空间组织，将气象特征在空间维聚合为每个时刻一个 memory token，仅保留时间位置表示；No delta 移除气象一阶差分分支，仅使用风场分支和背景场分支；Single branch 将风场、背景场和差分三路解耦编码改为对全部气象通道使用单一分支；No synoptic 移除背景场分支，并在差分分支中屏蔽背景场通道；No wind 移除风场分支，并在差分分支中屏蔽风场通道。

\casetable{表 3.}{消融实验结果。}{\input{../../paper_table_use/result/table03_ablation_results.tex}}

首先，去除功率上下文后的 Weather-only 模型 MAE 增加 4.27\%，RMSE 增加 2.94\%。将功率输入替换为随机值或零值后，MAE 分别增加 4.72\% 和 4.70\%，与 Weather-only 的退化幅度接近。这说明模型性能提升来自真实功率状态信息，而不是额外分支或参数量本身。将功率编码器替换为 Basic GCN 后，MAE 增加 2.29\%，RMSE 增加 1.96\%，表明图结构功率建模是有效的，但完整的多尺度功率动态编码能够取得更低误差。

其次，气象模块消融显示，气象空间记忆和气象变化建模对预测结果有直接影响。关闭气象空间记忆后，MAE 和 RMSE 分别增加 5.33\% 和 5.01\%；去除气象差分分支后，MAE 和 RMSE 分别增加 4.62\% 和 4.83\%；将解耦气象分支合并为单分支后，MAE 和 RMSE 分别增加 8.88\% 和 8.76\%。这些结果说明，相应气象编码组件被移除后，测试集误差均出现上升。

最后，气象通道消融表明风场信息对测试集误差影响最大。屏蔽背景气象变量后，MAE 和 RMSE 分别增加 11.66\% 和 10.95\%；屏蔽风场相关变量后，MAE 和 RMSE 分别增加 97.79\% 和 96.58\%，\(R^2\) 降至 0.3538。这说明在本文数据集上，风场相关变量是不可缺少的气象输入。

\casefigure{\textwidth}{../../paper_image_use/result/fig07_ablation_delta_vs_main.png}{图 7.}{消融模型相对于完整模型的误差增幅。}

\subsection{误差分布分析}

图 8 给出了不同模型在测试样本上的绝对误差分布。本文模型的误差分布整体低于其他基线，说明其优势不仅体现在平均指标上，也体现在测试样本层面的误差分布上。与此同时，各模型均存在较长的高误差尾部，表明在部分复杂天气或功率波动样本上，24 步集群总功率预测仍具有较高难度。

\casefigure{\textwidth}{../../paper_image_use/result/fig08_error_distribution.png}{图 8.}{不同模型的绝对误差分布。}

\subsection{不同预测步长上的结果}

为进一步分析模型在 24 步预测窗口内的表现，表 4 比较了本文模型和 Weather-only 在 H1、H6、H12、H24 以及全预测窗口平均意义下的误差。图 9 展示了两个模型在所有预测步长上的 MAE、RMSE 以及 MAE 差值。

\casetable{表 4.}{不同预测步长上的误差对比。}{\input{../../paper_table_use/result/table04_horizon_summary.tex}}

在首个预测步 H1 上，本文模型的 MAE 为 0.0430，显著低于 Weather-only 的 0.0747；RMSE 也从 0.0956 降至 0.0586。该结果说明近期功率上下文对预测窗口开始阶段最为有效，因为此时当前运行状态和短期惯性对目标功率仍具有较强约束作用。

随着预测步长增加，本文模型与 Weather-only 之间的差距逐渐缩小。在 H6 和 H12 上，本文模型仍保持较小优势；到 H24 时，Weather-only 的 MAE 和 RMSE 略低于本文模型。该结果表明，功率上下文带来的误差降低主要集中在预测窗口前段。尽管如此，在 1--24 步平均意义下，本文模型仍保持更低的 MAE 和 RMSE。

\casefigure{\textwidth}{../../paper_image_use/result/fig09_ours_weather_horizon_residual.png}{图 9.}{本文模型与 Weather-only 的逐步长误差及差值。}

\subsection{Discussion}

实验结果首先说明，区域集群总功率预测中未来气象预报仍是主导性信息来源。Weather-only 在所有基线中表现较强，表明 Pangu-Weather 网格化气象场已经能够刻画预测窗口内的主要外生驱动。然而，仅依赖气象场并不能完全描述目标窗口开始前的实际运行状态。本文模型在 Weather-only 的基础上进一步引入近期场站级功率，并将其作为动态上下文注入气象记忆，使预测过程不只是从未来气象场直接回归功率，而是在当前场群状态约束下读取未来气象信息。

逐步长结果揭示了这一条件化机制的作用边界。功率上下文带来的改善主要集中在预测窗口前段，尤其在 H1 上最为明显；随着预测步长增加，本文模型与 Weather-only 的差距逐渐缩小，H24 上 Weather-only 甚至略优。这一现象并不削弱功率上下文的意义，而是说明不同信息源在时间轴上的作用不同：目标窗口前最新功率更适合约束近端预测，因为其中包含当前出力水平、短期惯性和场站间同步变化；远端预测则更多受未来气象预报质量支配。因此，本文方法的核心不是用历史功率替代气象预报，而是用近期功率对气象驱动预测进行状态条件化。

消融实验进一步支持上述解释。随机功率和零功率输入接近 Weather-only，说明收益并非来自额外参数量或功率分支形式，而依赖真实功率状态本身。Basic GCN 的性能弱于完整模型，表明仅用单一时间尺度的图编码不足以充分刻画场站级功率动态；点级、片段级和全局尺度之间的交互有助于形成更有效的功率上下文。气象侧消融则表明，空间气象记忆、气象变化分支和变量解耦编码共同影响模型读取未来气象信息的能力。

气象通道消融还给出了较清晰的物理解释。屏蔽背景气象变量会造成误差上升，但屏蔽风场相关变量会导致性能大幅退化，说明本文任务中的主要可预测性仍来自未来风场。背景场变量的作用更接近对大尺度天气状态的补充约束，而不能替代与风电出力直接相关的风场信息。这一结果也表明，模型性能并非来自简单堆叠更多气象通道，而是依赖气象变量与风电出力之间的物理相关性。

从应用角度看，本文结果支持一种符合时间可用性约束的区域集群预测流程：在每天预测目标窗口开始前，使用最新可获得的场站级功率表征当前运行状态，同时使用 Pangu-Weather 预报提供未来气象驱动。该流程对预测窗口前段尤其有价值；对于更远预测步长，进一步结合气象预报不确定性有望提升模型对远端预测的适应能力。本文以吉林区域风电场群为案例开展系统验证，后续可在更多区域和不同装机结构的风电场群上检验模型迁移能力，并进一步面向极端爬坡、限电恢复和概率预测等调度相关场景扩展。
